\chapter{Unit 2: Empathy, Observation and Problem Identification}
\label{chap:unit2_qb}

\section*{Practice Question Set (30 MCQs)}
\addcontentsline{toc}{section}{Practice Question Set (30 MCQs)}

\mcqitem{1}{What is empathy in the context of Design Thinking?}{Feeling pity or sympathy for users}{The deep understanding of user feelings, motivations, behaviors, and unmet needs}{Selling products using emotional advertisements}{Guessing what users want based on designer intuition}

\mcqitem{2}{Which statement correctly differentiates a user "need" from a "solution"?}{"Users need an iPhone app" is a need; "Faster communication" is a solution}{"Users need timely bus arrival info" is a need; "A mobile GPS tracker" is a solution}{Needs and solutions are identical terms in design thinking}{A solution is always identified before the user need}

\mcqitem{3}{What is "Computational Empathy" in software engineering?}{Giving computers feelings and consciousness}{Designing systems that adapt gracefully to user context, errors, and emotional states}{Calculating mathematical probabilities of human emotions}{Replacing human customer support agents with chatbots}

\mcqitem{4}{Qualitative user research primarily aims to answer which question?}{"How many users clicked the button?"}{"What is the average transaction value?"}{"Why and how do users experience this problem?"}{"What percentage of users are on iOS?"}

\mcqitem{5}{What is "Data Triangulation" in user research?}{Surveying exactly three users}{Cross-verifying findings using multiple independent research methods and data sources}{Arranging interview transcripts in triangular shapes}{Using three different chart types in a presentation}

\mcqitem{6}{What does the "E" stand for in the AEIOU observation framework?}{Emotion}{Environment}{Evaluation}{Efficiency}

\mcqitem{7}{In the AEIOU framework, "Interactions" refers to:}{Building relationships between project managers}{Exchanges between person-to-person, person-to-device, or person-to-object}{Social media marketing campaigns}{Network server communications}

\mcqitem{8}{In user interviews, why are open-ended questions preferred over closed-ended questions?}{They elicit rich stories, past behaviors, and detailed context rather than simple "Yes/No"}{They take much less time to administer}{They are easier to grade mathematically}{They prevent participants from talking too much}

\mcqitem{9}{Which of the following is an example of a LEADING question that should be AVOIDED?}{"Tell me about the last time you booked a doctor's appointment."}{"What challenges did you face during your recent online order?"}{"Don't you agree that our new checkout screen is very fast and helpful?"}{"How did you feel when the payment failed?"}

\mcqitem{10}{The "Laddering" or "5 Whys" technique is primarily used to:}{Count the number of steps in a software flow}{Uncover deep underlying root causes and emotional drivers behind a surface behavior}{Grade user intelligence}{Prolong the duration of an interview session}

\mcqitem{11}{What is an "Explicit Need"?}{An unspoken subconscious desire}{A directly verbalized and clearly stated user requirement}{A requirement that violates government regulations}{An engineering specification hidden in code}

\mcqitem{12}{What is an "Implicit (Latent) Need"?}{An unstated expectation inferred through observing user workarounds, body language, and hesitation}{A feature written in the user manual}{A bug reported on an official feedback form}{A need that is impossible to satisfy}

\mcqitem{13}{When a user has to enter the exact same shipping address across 4 different checkout screens, what type of pain point is this?}{Financial pain point}{Time / Process pain point}{Aesthetic pain point}{Environmental pain point}

\mcqitem{14}{Hidden delivery charges revealed only on the final payment page create which type of user pain point?}{Process pain}{Financial pain}{Technical pain}{Physical pain}

\mcqitem{15}{A User Persona in Design Thinking should be created based on:}{A designer's personal assumptions and stereotypes}{Empirical data, user research patterns, behaviors, and motivations}{Purely demographic averages (age, gender, income) without behavioral context}{Fictional characters from popular movies}

\mcqitem{16}{What is the primary difference between a Primary Persona and a Secondary Persona?}{Primary personas are real people; secondary personas are robots}{Primary persona is the core target user whose needs dictate system design; secondary needs must not compromise primary goals}{Secondary personas pay for the software; primary personas only use it}{Primary personas are older adults; secondary personas are youth}

\mcqitem{17}{Which of the following is NOT one of the four main quadrants of a standard Empathy Map?}{Says}{Thinks}{Codes}{Feels}

\mcqitem{18}{In an Empathy Map, where do observable actions, habits, and physical workarounds belong?}{Thinks}{Feels}{Does}{Says}

\mcqitem{19}{Why is it important to distinguish between "Says" and "Does" in an Empathy Map?}{Users often state what is socially desirable but behave differently in reality}{"Says" is always accurate while "Does" is misleading}{Designers only care about what users say}{"Does" only applies to hardware devices}

\mcqitem{20}{What is a "Point-of-View" (POV) statement in the Define stage?}{A legal copyright statement}{A meaningful problem statement combining User, Need, and Insight}{The marketing slogan of the company}{A technical bug report}

\mcqitem{21}{How do design teams convert a POV problem statement into a launchpad for creative brainstorming?}{By writing software requirement specifications (SRS)}{By creating "How Might We" (HMW) questions}{By immediately purchasing off-the-shelf software}{By calculating revenue forecasts}

\mcqitem{22}{What makes a "How Might We" (HMW) question effective?}{It is so narrow that it includes the exact solution}{It is broad enough to allow creative exploration, yet narrow enough to provide clear focus}{It can be answered with a simple "Yes" or "No"}{It focuses exclusively on technology costs}

\mcqitem{23}{In the AEIOU framework, a handwritten cheat-sheet taped to a computer monitor is categorized as an:}{Activity}{Environment}{Object}{User}

\mcqitem{24}{During observational field research, what is the best practice for researchers?}{Immediately intervene and correct the user when they make an error}{Observe silently, timestamp events, and record actual behavior without judgment}{Pitch new product ideas to the user while they work}{Ignore the environment and only look at the computer screen}

\mcqitem{25}{What is the primary ethical responsibility of a design researcher during interviews?}{Guaranteeing that participants will receive free products}{Obtaining informed consent, ensuring data privacy, and allowing withdrawal at any time}{Recording audio without informing the user to ensure natural responses}{Convincing the participant that the current product is perfect}

\mcqitem{26}{Which of the following is a classic indicator of an unmet user need during observation?}{The user completes the task effortlessly in 5 seconds}{The user creates an improvised workaround (e.g., using personal spreadsheets or paper notes)}{The user follows the official documentation without hesitation}{The user smiles and praises the software}

\mcqitem{27}{How does quantitative research support qualitative research in the empathy phase?}{It replaces qualitative research completely}{It validates how widespread a discovered qualitative pain point is across a broader population}{It generates emotional empathy maps}{It eliminates the need to interview real users}

\mcqitem{28}{An actionable design insight is best described as:}{A raw user quote copied directly into a report}{A statistical graph showing website traffic}{A deep, non-obvious realization that connects user behavior to underlying human motivations}{A list of software features requested by management}

\mcqitem{29}{What type of pain point occurs when a user feels anxious and embarrassed because an ATM display exposes their account balance to everyone in line?}{Physical pain}{Emotional / Privacy pain}{Financial pain}{Operational pain}

\mcqitem{30}{Which of the following defines a Target Archetype?}{A strict mathematical demographic model}{A generalized behavioral profile capturing shared goals, mental models, and usage contexts}{A 3D CAD model of an average human body}{The software architecture blueprint}


\newpage
\section*{Answer Key \& Step-by-Step Explanations}
\addcontentsline{toc}{section}{Answer Key \& Explanations}

\begin{answerkeytable}{Unit 2: Quick Answer Key}
\begin{tabular}{c c c c c c c c c c}
\toprule
\textbf{Q1} & \textbf{Q2} & \textbf{Q3} & \textbf{Q4} & \textbf{Q5} & \textbf{Q6} & \textbf{Q7} & \textbf{Q8} & \textbf{Q9} & \textbf{Q10} \\
\textbf{B} & \textbf{B} & \textbf{B} & \textbf{C} & \textbf{B} & \textbf{B} & \textbf{B} & \textbf{A} & \textbf{C} & \textbf{B} \\
\midrule
\textbf{Q11} & \textbf{Q12} & \textbf{Q13} & \textbf{Q14} & \textbf{Q15} & \textbf{Q16} & \textbf{Q17} & \textbf{Q18} & \textbf{Q19} & \textbf{Q20} \\
\textbf{B} & \textbf{A} & \textbf{B} & \textbf{B} & \textbf{B} & \textbf{B} & \textbf{C} & \textbf{C} & \textbf{A} & \textbf{B} \\
\midrule
\textbf{Q21} & \textbf{Q22} & \textbf{Q23} & \textbf{Q24} & \textbf{Q25} & \textbf{Q26} & \textbf{Q27} & \textbf{Q28} & \textbf{Q29} & \textbf{Q30} \\
\textbf{B} & \textbf{B} & \textbf{C} & \textbf{B} & \textbf{B} & \textbf{B} & \textbf{B} & \textbf{C} & \textbf{B} & \textbf{B} \\
\bottomrule
\end{tabular}
\end{answerkeytable}

\begin{expbox}[Detailed Pedagogical Explanations]
\begin{enumerate}[leftmargin=6mm, itemsep=2.5mm]
  \item \textbf{[Q1: Option B]} Empathy is understanding the user's emotions, experiences, and implicit needs from their perspective.
  \item \textbf{[Q2: Option B]} Needs are underlying goals/outcomes; solutions are specific implementations.
  \item \textbf{[Q3: Option B]} Computational empathy designs software that adapts to human context, cognitive load, and emotions.
  \item \textbf{[Q4: Option C]} Qualitative research explores the rich "Why" and "How" behind human behavior.
  \item \textbf{[Q5: Option B]} Triangulation compares multiple distinct research methods to validate findings.
  \item \textbf{[Q6: Option B]} In AEIOU: Activities, Environments, Interactions, Objects, Users.
  \item \textbf{[Q7: Option B]} Interactions cover exchanges between humans, devices, and physical objects.
  \item \textbf{[Q8: Option A]} Open-ended questions allow participants to share realistic stories and unprompted context.
  \item \textbf{[Q9: Option C]} Leading questions nudge the participant toward a desired positive answer and introduce heavy bias.
  \item \textbf{[Q10: Option B]} Laddering (5 Whys) drills down from superficial actions to root motivations.
  \item \textbf{[Q11: Option B]} Explicit needs are directly stated and clearly communicated by users.
  \item \textbf{[Q12: Option A]} Implicit needs are unstated desires uncovered through observations and workarounds.
  \item \textbf{[Q13: Option B]} Re-entering identical information is a classic time and workflow process friction.
  \item \textbf{[Q14: Option B]} Unexpected costs and hidden pricing friction represent financial pain points.
  \item \textbf{[Q15: Option B]} Personas must be grounded in real behavioral research data, not designer stereotypes.
  \item \textbf{[Q16: Option B]} Primary personas represent the core user whose needs must be satisfied first without compromise.
  \item \textbf{[Q17: Option C]} Empathy map quadrants are Says, Thinks, Does, and Feels (plus Pains \& Gains).
  \item \textbf{[Q18: Option C]} "Does" captures observable physical actions, workflows, and behavioral adaptations.
  \item \textbf{[Q19: Option A]} People often say what sounds good or polite, but their actual behavior reveals true needs.
  \item \textbf{[Q20: Option B]} A POV statement synthesizes User + Need + Surprising Insight.
  \item \textbf{[Q21: Option B]} "How Might We" questions reframe problem insights into constructive ideation prompts.
  \item \textbf{[Q22: Option B]} Good HMWs strike a balance: broad enough for variety, narrow enough for actionable focus.
  \item \textbf{[Q23: Option C]} Physical artefacts, cheat-sheets, and improvised tools fall under 'Objects'.
  \item \textbf{[Q24: Option B]} Neutral observation without interrupting or correcting yields authentic behavioral data.
  \item \textbf{[Q25: Option B]} Respecting privacy, gaining informed consent, and protecting participants are ethical imperatives.
  \item \textbf{[Q26: Option B]} Workarounds signal that existing official tools fail to satisfy a crucial user need.
  \item \textbf{[Q27: Option B]} Quantitative data tests the scale and frequency of qualitative insights.
  \item \textbf{[Q28: Option C]} Insights go beyond surface facts to explain the underlying "Why".
  \item \textbf{[Q29: Option B]} Anxiety, vulnerability, and embarrassment are emotional/psychological pain points.
  \item \textbf{[Q30: Option B]} Archetypes represent behavioral patterns and shared mental models across user groups.
\end{enumerate}
\end{expbox}
